% =====================================================================
%  GMN3 §10.4--10.8  $A_1$ 理论的例子（续）
%  （$N_f=1$ / 其他 $N_f$ / $N_f=4$ / $\CN=2^*$ / millipede 展开）
%  原文：Lamination-Paper.tex 5655--5863 行；PDF p.93--98（印刷 p.92--97）
%  编号公式 (10.35)--(10.48) 共 14 条（equation 5 处 + align 6 行 + align 3 行）；
%  插图 9 幅（nf1-maximal / nf1 / nf1-l0 / nf2-l / nf3-l / nf4 / nf4-lstar /
%  n2star-triang / millipede）
%  注：BRIEF 分块表记本块 25 条公式，与原 PDF 不符；§10 全节 48 条的其余 34 条
%      落在 ch12（§10.1--10.3），本块末条 (10.48) 正好收满 §10，故以 14 条为准。
%  跨块引用（节号、式号）暂写原文编号并加 % TODO，合稿时统一改 \ref/\eqref
% =====================================================================

\subsection{$SU(2)$，$N_f = 1$}

下面我们处理带一重基础表示味的 $SU(2)$ 规范理论。正如文献 \cite{Gaiotto:2009hg} 第 10.2 节所解释的，这是一个 $A_1$ 理论，其中我们取 $C = \IC\IP^1$，带两个不规则奇点 $\CP_1$ 与 $\CP_2$，它们分别承载 $\phi_2$ 的 $3$ 阶极点与 $4$ 阶极点。拓扑上，我们在每个奇点周围挖去一个圆盘，并在 $S^1(\CP_1)$ 上标出一个点、在 $S^1(\CP_2)$ 上标出两个点。

该理论中的一个叶层样本见图 \ref{fig:nf1-maximal}。整数权重 $k_i$ 受如下限制：
\begin{equation} \label{eq:nf1-cond}
2 k_1 + k_4 + k_5 + k_6 = 0, \quad k_2 + k_3 + k_4 = 0, \quad k_2 + k_3 + k_5 + k_6 = 0, \quad k_4 \ge 0, \quad k_5 \ge 0, \quad k_6 \ge 0.
\end{equation}
条件 \eqref{eq:nf1-cond} 的解由 $k_4, k_5, k_6$ 决定：它们都 $\ge 0$，且 $k_4 + k_5 + k_6$ 为偶数。
\insfig{nf1-maximal}{$N_f=1$ 的 $SU(2)$ 理论中的一族叶层。}
于是我们得到了一个由这样的 $(k_4, k_5, k_6)$ 参数化的叶层族。

然而，这些并不是全部可能的叶层，因为我们所选的曲线并非最一般的可能选法。与 $N_f = 0$ 理论一样，我们可以通过一种 Dehn 扭转生成更多的叶层：把 $S^1(\CP_2)$ 的一个邻域连续旋转角度 $\pi$（于是两个标出的点互换）。这一扭转把路径 $\wp_4$、$\wp_5$、$\wp_6$ 拖成新的路径。对我们上面构造的叶层反复作顺时针或逆时针扭转，就能得到几乎所有可能的叶层。与前面一样，我们用 $L_*$ 记由单条闭曲线组成的叶层，它不能由上述构造得到。

现在如下选取一个三角剖分。边界圆 $S^1(\CP_2)$ 被分成两段。其中一段只与一条边相接，我们把那条边标为 $X$。另一段与两条边相接，我们把它们标为 $Y$ 与 $Z$。这三条边的另一端都落在边界圆 $S^1(\CP_1)$ 上。沿一条起止都在 $S^1(\CP_1)$ 上的弧顺时针走，我们会依次遇到 $Z$、$X$、$Y$。见图 \ref{fig:nf1}。
\insfig{nf1}{$N_f = 1$ 的 $SU(2)$ 理论中的一个三角剖分。}

设 $L_0$ 为图 \ref{fig:nf1-l0} 所示的叶层。
\insfig{nf1-l0}{$N_f = 1$ 的 $SU(2)$ 理论中的叶层 $L_0$。}
其余基本叶层 $L_n$ 由 $L_0$ 绕 $\CP_2$ 作上述扭转得到。计算它们的 vev 便给出，例如
\begin{align}
 \inprod{L_{-2}}' &= Y + XY + YZ + X Y Z^2 + 2 XYZ, \\
 \inprod{L_{-1}}' &= X + XZ, \\
 \inprod{L_{0}}' &= Z, \\
 \inprod{L_{1}}' &= 1/Y, \\
 \inprod{L_{2}}' &= 1/X + 1/XY, \\
 \inprod{L_3}' &= 1/Z + 1/XZ + 1/X Y^2 Z + 1/YZ + 2/XYZ.
\end{align}
一般地，$\inprod{L_k}'$ 与 $\inprod{L_{1-k}}'$ 由同时互换 $Z \leftrightarrow 1/Y$、$X \leftrightarrow 1/X$ 相联系。

我们还有
\begin{equation}
\inprod{L_*}' = \frac{1 + X + XY + XYZ}{\sqrt{XYZ}}.
\end{equation}

\subsection{$SU(2)$，其他 $N_f$}

对所有的 $SU(2)$、$N_f > 0$ 理论，都可以类似地描述其叶层。这里我们只局限于几条简短的评注。

$N_f=2$ 的理论有两种实现，见文献 \cite{Gaiotto:2009hg} 第 10.3 节与第 10.5 节。我们考虑它的「第一种实现」：其中 $C$ 为 $\IC\IP^1$，带两个不规则奇点，每个奇点承载两条斯托克斯射线。一个叶层样本见图 \ref{fig:nf2-l}。这个叶层在通常意义下是极大的，即不能再向它添加曲线。若改变附着在曲线上的权重，我们就能得到更一般的叶层族。
\insfig{nf2-l}{$N_f=2$ 的 $SU(2)$ 理论在其第一种实现下的一个叶层样本。}

$N_f = 3$ 的理论通过取 $C$ 为 $\IC\IP^1$ 得到：它带一个不规则奇点（承载两条斯托克斯射线）与两个正则奇点。一个叶层样本见图 \ref{fig:nf3-l}。
\insfig{nf3-l}{$N_f=3$ 的 $SU(2)$ 理论中的一个叶层样本。}

\subsection{$SU(2)$，$N_f = 4$}

接下来我们简要考虑 $N_f = 4$ 的 $SU(2)$ 理论。该理论对应 $C = \IC\IP^1$ 带 $4$ 个正则奇点的情形。对 $(\zeta, u)$ 的某些取法，$T_\WKB(\zeta, u)$ 就是图 \ref{fig:nf4} 所画的「四面体」三角剖分。
\insfig{nf4}{$\IC\IP^1$ 的「四面体」三角剖分，它出现在 $N_f=4$ 的 $SU(2)$ 理论中。顶点 $i$ 与顶点 $j$ 之间的边标为 $X_{ij}$。}
然后考虑图 \ref{fig:nf4-lstar} 所示的叶层 $L_*$。应用文献 \cite{Drukker:2009tz} 的规则可知，对应的算子 $L_*(\zeta)$ 是某一对偶标架下的 Wilson 圈。
\insfig{nf4-lstar}{$N_f=4$ 的 $SU(2)$ 理论中的叶层 $L_*$。}
利用 traffic rule，它的真空期望值为
%
\begin{equation}
\inprod{L_*} = {\Tr} \, L M_{23} R M_{13} L M_{14} R M_{24}
\end{equation}
%
计算出来是
%
\begin{equation}\label{eq:Nf4-Wilson}
\inprod{L_*} = \frac{1 + X_{13} + X_{24} + X_{13} X_{24} + X_{13} X_{14} X_{24} + X_{13} X_{23} X_{24} + X_{13} X_{14} X_{23} X_{24}}{\sqrt{X_{13} X_{14} X_{23} X_{24}}}.
\end{equation}
%
于是这个线算子支持 $7$ 个 BPS 态，它们的电荷各不相同。

\subsection{$SU(2)$，$\CN=2^*$}\label{subsec:SU2-N2star-Examples}

\insfig{n2star-triang}{$\CN=2^*$ 理论的通用 WKB 三角剖分。}

最后我们考虑最有趣的例子：$\CN=2^*$ 理论。正如我们在第 5.4 节% TODO \ref{subsec:Formal-N2-star}
所回顾的，该理论对应取 $C$ 为带一个正则奇点的单穿刺环面。相关的二次微分为 $\phi_2 = u + m^2 \wp(z|\tau)$。这里有两个转向点，且所有 WKB 三角剖分的拓扑都相同：两个三角形，把环面的某个基本区域一分为二。不失一般性，我们可以用三条直线 $(0,1)$、$(0,\tau)$、$(0,\tau+1)$ 表示这些边，如图 \ref{fig:n2star-triang} 所示。我们把相应的交比记作 $X$、$Y$、$Z$。于是绕穿刺点的单值为 $\mu = - XYZ$。绕环面的基本闭链 $z \to z+1$ 与 $z \to z+\tau$ 的和乐 ${\rm Tr}(A)$ 与 ${\rm Tr}(B)$ 很容易算出：这两条闭链各自穿过两条边，因而具有标准的三项展开。对闭链 $z \to z+1+\tau$ 也一样，它给出 ${\rm Tr}(AB)$。于是我们有展开
%
\begin{align}\label{eq:N2-star-traffic}
\inprod{L_{0,1}}' &= {\rm Tr}(A) = \sqrt{Z Y} + \sqrt{Z/Y} + 1/\sqrt{Z Y}, \\
\inprod{L_{1,0}}' &= {\rm Tr}(B) = \sqrt{X Z} + \sqrt{X/Z} + 1/\sqrt{X Z}, \\
\inprod{L_{1,1}}' &= {\rm Tr}(AB) = \sqrt{Y X} + \sqrt{Y/X} + 1/\sqrt{Y X}.
\end{align}
%
我们可以把这些与形式 Wilson 线算子 (5.24)% TODO \eqref{eq:star-formal-Wilson}
作比较：于是 $W_1$ 对应 ${\rm Tr}(A)$ 的（经典极限），依此类推，与第 8.4 节% TODO \ref{subsec:N2-star-moduli}
的情形完全相同。作为检验，由 \eqref{eq:N2-star-traffic} 可知：若把 $X, Y, Z$ 看成交换变量，则代入 \eqref{eq:N2-star-traffic} 确实给出
%
\begin{equation}
({\rm Tr}(A))^2 + ({\rm Tr}(B))^2 + ({\rm Tr}(AB))^2 - {\rm Tr}(A) {\rm Tr}(B) {\rm Tr}(AB) = \mu + 2 + \frac{1}{\mu}
\end{equation}
%
其中我们用了 $\mu = - XYZ$。这与方程 (8.9)% TODO \eqref{eq:moduli-equation}
相符，并且是 (5.34)% TODO \eqref{eq:quant-mod-eq}
的经典极限。

利用「提升原理」（promotion principle），由这些表达式可知：第 5.4 节% TODO \ref{subsec:Formal-N2-star}
构造的形式线算子确实计算 $\CN=2^*$ 理论的保护自旋特征标。

\subsection{millipede 展开}

在 $A_1$ 理论中，「原味」BPS 态有一个简单的描述：它们对应 $(2,0)$ 理论在 $C$ 上的、总质量有限的 BPS 弦。这归结为或者闭合、或者起止于二次微分 $\phi_2$ 之零点的 WKB 曲线。这一图像由文献 \cite{Klemm:1996bj,Shapere:1999xr} 引入，并在文献 \cite{Gaiotto:2009hg} 中起了重要作用。本文引入了带框 BPS 态，自然要问：它们是否也由类似的几何图像刻画。我们相信答案是肯定的：本节简要勾勒一个方案，把带框 BPS 简并度 $\fro(L, \gamma, \zeta)$ 直接等同于计数某些对象，我们称之为\emph{以 $L$ 为躯干、相位为 $\arg \zeta$ 的 millipede}（千足虫）。

我们这里不一般地定义这些对象，而只限于叶层 $L$ 就是一条闭圈、并承载 $SL(2,\IC)$ 之基础表示的情形。此时一个 millipede 是 $\Sigma$ 上的一条闭的有向曲线 $\xi$，它可以如下分成若干段。其中一些段称为\emph{躯干段}（body segment）：每条躯干段都是 $L$ 的相应一段到 $\Sigma$ 的提升。我们要求 $L$ 的每个点恰好提升为躯干段上的一个点：于是这些躯干段构成 $L$ 到 $\Sigma$ 的一种断裂提升。其余的段称为\emph{腿}（leg）。腿是沿之 $\lambda/\zeta$ 为实数、且取非负定向的有向曲线。（这样的段是 $C$ 上 WKB 曲线到 $\Sigma$ 的提升。）millipede 的\emph{荷}定义为 $\xi$ 在 $H_1(\Sigma, \IZ)$ 中的类。我们提出：带框 BPS 简并度应当「计数」以 $L$ 为躯干、相位为 $\arg \zeta$、荷为 $\gamma$ 的 millipede $\xi$。

我们来看构造 millipede 的最基本方式。从 $L$ 的一个断裂提升出发。在提升从 $\Sigma$ 的一叶跳到另一叶的地方，它带有边界点 $p_1$、$p_2$（它们投影到 $C$ 的同一个点 $p$）。于是 $\xi$ 的躯干不闭合。为修补这一点，我们在每个边界点 $p$ 处接上一条腿段：它从 $p_1$ 沿 $\Sigma$ 走到 $\lambda$ 的一个零点，再沿 $\Sigma$ 的另一叶返回到 $p_2$。（注意，只有当存在一条从 $p$ 通向某个转向点的 WKB 曲线时这才可能；而这类「特殊 WKB 曲线」只有有限多条，所以对给定的 $L$，接腿的方式只有有限多种。）把这一图像投影到 $C$，就解释了这一术语——见图 \ref{fig:millipede}。
\insfig{millipede}{一个 millipede 到 $C$ 的投影的一部分。躯干画成紫色，腿画成黑色。这个特定的 millipede 属于最简单的一类：所有腿都终止在转向点上，转向点画成橙色的叉。}

与「原味」BPS 态的情形一样，计数这些对象的精确规则有些微妙。最简单的情形就是我们刚讨论的：所有腿都从 $L$ 通向转向点。这样的 millipede 是孤立的（没有模）。如果只有这类 millipede，并且每个 millipede 在每个转向点上都只有一条腿终止，那么用 traffic rule 就比较容易看出：以 $L$ 为躯干、荷为 $\gamma$ 的 millipede 的数目，确实等于数目 $\fro(L, \gamma)$。

不过，更一般地，可能有若干条腿终止在同一个转向点上；在这些情形，该 millipede 有时贡献 $-1$ 而不是 $+1$。此外，还可能存在一些腿的起点与终点都在 $L$ 上的 millipede。这样的 millipede 有模（腿可以沿 $L$ 来回滑动）。我们猜想：孤立的 millipede 贡献自旋为零的带框 BPS 态，非孤立的 millipede 贡献更高自旋的带框 BPS 多重态；至少在一些简单例子中是这样，但我们尚未系统地研究这一问题。
